\section{Introduction}

\subsection{Why Are We Using Servers?}

Most scientific models and datasets are too large to run efficiently on a laptop.
The Fluids Lab servers, \texttt{Pelee} and \texttt{Rondeau}, provide access to faster CPUs, GPUs, larger memory, and shared storage.

You may think of them as powerful Linux computers that you access remotely.

\begin{center}
Your Laptop $\longleftrightarrow$ Pelee / Rondeau
\end{center}

You work from your laptop, but the calculations happen on the server.

\subsection{Important Concepts}

\paragraph{Multi-user systems.}
Many researchers share the same server.

\paragraph{CPU versus GPU.}
A CPU is designed for general-purpose computing. A GPU is designed for highly parallel computations.

\paragraph{Home directory.}
Every user has a home directory:

\begin{lstlisting}
/home/username
\end{lstlisting}
Note: for historical reasons at UWaterloo Math, $/home$ is an alias for $/u$, although they have the same purpose and meaning.

\paragraph{Persistence.}
Files remain on the server after you log out.

\subsection{Learning Objectives}

\begin{itemize}
\item connect to Pelee and Rondeau;
\item navigate the Linux filesystem;
\item transfer files;
\item run Python and Julia programs;
\item monitor resource usage.
\end{itemize}
